% Canon: SGA-SRC-001254.
\par\medskip\noindent
1.1.5.  Recíprocamente, supuestos dados los grupos $P$ y $G$ de
$\underline{T}$, supongamos dado además un bitorsor

\begin{equation}\tag{1.1.5.1}
j \; : \; E \longrightarrow P
\end{equation}
bajo $(G_{P},G_{P})$ (NB: las dos acciones de $G_{P}$ sobre $E$, o,
lo que es lo mismo, de $G$ sobre $\sisi{P}{E}$, se sobreentienden en las
notaciones), y, por otra parte, un sistema de isomorfismos de bitorsors
sobre una base variable $S$,

\begin{equation}\tag{1.1.5.2}
\varphi_{p,p'} \; : \; E_{p}E_{p'} \;\overset{\sim}{\longrightarrow}\; E_{pp'} \qquad (p,p' \in P(S)),
\end{equation}
compatibles con el cambio de base $S' \longrightarrow S$; en otras
palabras, un isomorfismo de bitorsors (1.1.3.2). Supongamos que los
diagramas correspondientes (1.1.4.1) son conmutativos. Afirmo que
\underline{existe sobre} $E$ \underline{una única estructura de
extensión de} $P$ \underline{por} $G$, \underline{tal que los datos
precedentes se deducen de esta estructura mediante las construcciones
de} 1.1.2 \underline{y} 1.1.3.

Observemos primero que, al tomar en (1.1.5.2) como $p,p'$ la sección
unidad $1$ de $P$, se obtiene un isomorfismo

\begin{equation}\tag{1.1.5.3}
\varphi_{1,1} \; : \; E_{1} \;\overset{\sim}{\longleftarrow}\; E_{1} \overset{G}{\times} E_{1} \qquad ,
\end{equation}
de donde, tomando el producto contraído (a izquierda, digamos) de ambos
miembros con el bitorsor inverso $E_{1}^{-1}$ de $E_{1}$ [2], se obtiene
un isomorfismo canónico

\begin{equation}\tag{1.1.5.4}
E_{1} \;\simeq\; \text{bitorsor unidad } {}_{s}G_{d}
\end{equation}
($=G$ actuando sobre sí mismo mediante traslaciones a izquierda y a
derecha), teniendo en cuenta
%%%%%%%%%%  scan idx 152  |  folio 141  %%%%%%%%%%
el isomorfismo canónico

\begin{equation*}
E_1^{-1}.E_1 \;\; \simeq \;\; \text{bitorsor unidad } {}_sG_d
\end{equation*}
de loc.cit. (válido para todo bitorsor). Esto equivale a decir que se
tiene una sección $1_E$ de $E_1$, y por tanto de $E$:

\begin{equation}\tag{1.1.5.5}
1_E \in \Gamma(E_1) = \mathrm{Hom}(e,E_1) \subset \Gamma(E)
\end{equation}
(a saber, la imagen por (1.1.5.4) de la sección marcada de
${}_sG_d$), que satisface la relación

\begin{equation*}
g.1_E = 1_E.g \qquad (g \in G(S)) \quad ,
\end{equation*}
y, por construcción, esta sección está caracterizada por la condición

\begin{equation}\tag{1.1.5.6}
\varphi_{1,1}(1_E,1_E) = 1_E \qquad .
\end{equation}
Por otra parte, el dato de un morfismo de torsors $\varphi$ (1.1.3.2)
equivale manifiestamente al dato de un morfismo

\begin{equation}\tag{1.1.5.7}
\pi_E : E\times E \longrightarrow E \quad , \qquad \text{denotado también } (x,y) \longmapsto xy
\end{equation}
que satisface las condiciones

\begin{equation}\tag{1.1.5.8}
j(xy) = j(x)j(y) \qquad (x,y \in E(S))
\end{equation}
(es decir, $j$ es compatible con las estructuras multiplicativas
$\pi_E$, resp. $\pi$, sobre $E$, resp. $P$), y las tres relaciones
(1.1.3.0) (comenzando preferentemente por la última de ellas). Por
último, se define un morfismo de objetos de \underline{$T$}

\begin{equation*}
i : G \longrightarrow E
\end{equation*}
por la fórmula

\begin{equation}\tag{1.1.5.9}
i(g) = g.1_E = 1_E.g 
\end{equation}

%%%%%%%%%%  scan idx 153  |  folio 142  %%%%%%%%%%

(la identidad de los dos últimos miembros de esta fórmula ya se había
observado en la construcción de $1_E$). Sentado esto, afirmo que
\underline{la ley de multiplicación} $\pi_E$ \underline{sobre} $E$
\underline{hace de} $E$ \underline{un grupo que admite} $1_E$
\underline{como sección unidad, y que}
$j : E \rightarrow P$ \underline{e} $i : G \rightarrow E$
\underline{son homomorfismos de grupos que hacen del grupo} $E$
\underline{una extensión de} $P$ \underline{por} $G$.

En efecto, la asociatividad de la ley de composición de $E$ no es sino
la conmutatividad de los diagramas (1.1.4.1). La construcción de
$\pi_E$ mediante los isomorfismos $\varphi_{p,p'}$ (1.1.3.1) muestra,
por otra parte, que para todo $x \in E(S)$ las traslaciones a izquierda
y a derecha por $x$ en $E(S)$ son inyectivas. Como
$1_E1_E = 1_E$ por (1.1.5.6), se concluye de inmediato que $1_E$ es una
unidad bilateral de $E$. Dejamos al lector la comprobación de la
existencia de un inverso de un $x \in E(S)$ sobre $p \in P(S)$,
utilizando un isomorfismo canónico deducido de
$\varphi_{p,p^{-1}}$:

\begin{equation}\tag{1.1.5.10}
\sisi{P_{p^{-1}}}{E_{p^{-1}}} \;\simeq\;
(\sisi{P_p}{E_p})^{o} \qquad .
\end{equation}
Por tanto, $E$ es en efecto un \underline{grupo}. Ya se ha señalado que
$j$ es compatible con las leyes de composición (1.1.5.8), y lo mismo
ocurre con $i$, pues, para $g,g' \in G(S)$, se tendrá:

\begin{multline*}
i(g)i(g') = (g.1_E)(g'.1_E) = (g.1_E.g')\ 1_E \quad (\text{última relación} \\ \text{de } (1.1.3.0)) = \ g.1_E.g' \quad (\text{pues } 1_E \text{ es unidad}) = g.(1_E.g') = g.(g'.1_E) \\ (1.1.5.9) = (gg').1_E = i(gg'). \qquad .
\end{multline*}
Como $j$ es la proyección estructural de un torsor, es un epimorfismo
de objetos de \underline{$T$}; queda por ver que $i$ induce un
isomorfismo de $G$ con $Ker j = E_1$, lo cual está claro por la
definición (1.1.5.9) y la construcción de $1_E$.


%%%%%%%%%%  scan idx 154  |  folio 143  %%%%%%%%%%

1.1.6.  Admitiremos que los razonamientos precedentes demuestran que,
para un topos fijo \underline{$T$}, las construcciones anteriores
proporcionan una \underline{equivalencia entre la categoría de las
extensiones de grupos de} \underline{$T$} (1.1.1) (donde $P$ y $G$,
así como $E$, se suponen variables), \underline{y la categoría de los
sistemas} $(P,G,E,\varphi)$, \underline{donde} $P$ \underline{y} $G$
\underline{son grupos de} \underline{$T$}, $E$ \underline{un bitorsor
bajo} $(G_P,G_P)$ \underline{sobre} $P$, \underline{y, finalmente,}
$\varphi$ \underline{un isomorfismo de bitorsors} (1.1.3.2) (o, si se
prefiere, un sistema de isomorfismos de bitorsors (1.1.3.1) compatibles
con los cambios de base), \underline{que satisface la condición de
asociatividad expresada por la conmutatividad del diagrama} (1.1.4.2)
(o, si se prefiere, de los diagramas (1.1.4.1)).

Dejamos al lector el cuidado de explicitar de manera evidente, en
términos del diccionario precedente, la noción de imagen inversa (resp.
directa) de una extensión (1.1.1) por un morfismo de grupos

\begin{equation*}
P' \longrightarrow P \qquad (\text{resp. } G \longrightarrow G'\ ) \qquad ,
\end{equation*}
o por estas dos operaciones simultáneamente. Le dejamos también el
cuidado de explicitar una compatibilidad de la equivalencia de
categorías precedente con el funtor imagen inversa $f^*$ por un morfismo de
topos

\begin{equation*}
f : \underline{T}' \longrightarrow \underline{T} \qquad ;
\end{equation*}
para ello, conviene utilizar el dato $\varphi$ en la forma (1.1.3.2)
(con preferencia a los $\varphi_{p,p'}$ de (1.1.3.1)), que define
entonces un dato correspondiente $f^*(\varphi)$ sobre el topos
$\underline{T}'$.
